\section{System and Method}\label{sec:method}
\textbf{Simulator and circuits.} All boards run the same C++ statevector code: the state is $2^n$ complex float32 amplitudes (\AmpBytes\,B each) and each gate kernel (H, CNOT, controlled-phase) sweeps all $2^n$ indices, so the work of a circuit is (gate calls)$\times 2^n$. The four circuits are \hall\ ($n$ H gates), \bell\ (one H and $n{-}1$ CNOT), \qft\ ($n$ H and $n(n{-}1)/2$ controlled-phase gates) and \grover\ ($\lfloor\frac{\pi}{4}\sqrt{2^n}\rfloor$ iterations of a phase oracle on $|1\ldots1\rangle$ and a diffusion step, both built from recursively decomposed multi-controlled phase gates). Gate counts are derived from the circuit header and cross-checked by running the unmodified header with counting kernels (\GateCrossN\ of \GateCrossAll\ $(\text{circuit},n)$ pairs equal). The firmware and the PC reference build the same gate sequence, tested gate by gate for all four circuits up to $n=\GateSeqMax$ and for \grover\ at $n=\GateSeqGrover$ (the iteration count agrees for every $n$ up to \KMaxN).

\textbf{Boards.} Arduino \uno\ (\FcpuUno\,MHz, \MemSramUno\,B SRAM, static array), \espA\ NodeMCU (\FcpuEspLo\,MHz) and \espB\ DevKit (\FcpuEspHi\,MHz), one unit each. Boards return timing only; they do not return the final state.

\textbf{Phases.} \emph{Phase~1} runs the float32 build on the PC against Qiskit Aer (infidelity and total variation distance). \emph{Phase~2} runs the circuit locally on the board for $n=2,3,\ldots$ until a memory wall or the per-point time cap (\PointCapH\,h); each point is up to \Reps\ repetitions (fewer for long \grover\ points), reported as median and interquartile range (IQR) of the per-repetition time measured with the board's \texttt{micros()} around the circuit call (resolution \UnoResUs\us\ on the Uno, \EspResUs\us\ on the ESP boards). A wall is recorded only with positive evidence: a linker RAM overflow (Uno) or a failed allocation (ESP boards). A separate ESP32 timestamps the rising and falling edge of a pin raised around the same call (32-bit capture at the \ApbMHz\,MHz clock, \CounterTickNs\,ns per tick, WiFi/TCP link to the PC) and serves as an independent reference for \texttt{micros()}. The counter is \emph{uncalibrated}: a loopback self-test passed \SelfTestPass\ of \SelfTestAll\ scenarios with a constant offset of \SelfTestLo--\SelfTestHi\us, but one crystal generates and measures. \emph{Phase~3} offloads: the board sends a \ReqBytes-byte request (start byte, $n$, circuit, CRC-8) over serial at \BaudSp\,baud; the PC simulates the circuit with Qiskit Aer, draws one measurement outcome and answers with a \RespBytes-byte frame (outcome, CRC-8). The board times the round trip with \texttt{micros()}. One shot per request is sampled and compared with no reference; the round trip is a latency, not a correctness result.

\textbf{Offload decomposition.} A Phase~3 point is split into wire time ($(\ReqBytes{+}\RespBytes)$ bytes $\times$ \BitsPerByte\ bit / baud $=\WireMs$\,ms), the PC simulation time logged per point, and a residual (board compute, USB-serial latency, host scheduling); the three are medians of one point and do not add exactly.

\textbf{Scope and session selection.} We use $n\ge2$ and, in Phase~3, $n$ up to the board's Phase~2 wall. Several points were measured in more than one session (original and repeats). A single selector fixes which session is primary evidence; the other sessions feed only the reproducibility analysis. The selection was made late in the analysis, \emph{after} the data and the first scores were known (Section~\ref{sec:limits}). The audit found \MixedPoints\ points with mixed firmware hashes in the selected data; firmware versions are classified by their effect on the timed region (class A: reference; B: timing-relevant difference, as for all selected \espA\ \grover\ points, which ran a timer interrupt inside the timed region; X: unknown uncommitted version; hash-flagged points are not fitted).

\textbf{Weak points.} A point with fewer than \WeakMin\ \emph{independent} repetitions (Phase~3: requests) is \emph{weak}: it is never fitted, never in the primary error, and marked as such in the figures (hollow markers in Fig.~\ref{fig:f4}; grey in Fig.~\ref{fig:f2}, where hollow instead means held out; Fig.~\ref{fig:f1} draws lines only). Two \espB\ \grover\ points contain exactly repeated rows (\DupNineIndep\ and \DupTenIndep\ independent of \DupNineRows\ recorded rows); the cause is unverified and only independent rows are used.

\textbf{Counter validity.} A counter value is used only if the join produced pairs (R1), the number of pairs equals the number of repetitions (R2), no paired pulse carries a counter-side flag (R3), and the median is below the representable limit $T_{\max}$ = \TMaxS\,s (R4). Rule R5 rejects a pair whose counter/\texttt{micros()} ratio lies outside $[\RFiveLo,\RFiveHi]$ or whose counter value is within \RFiveMs\,ms of a multiple of \WrapS\,s (the $2^{32}$-tick wrap artefact). \emph{R5 was added after the data had been seen}; the agreement statistics of Section~\ref{sec:res-counter} are therefore conditional on this coarse agreement filter.

\textbf{Memory model.} A point fits if overhead $+\ \AmpBytes\cdot2^n\le B$, so $n_{\max}=\lfloor\log_2(B/\AmpBytes)\rfloor$, with $B$ read from low-$n$ runs only: the Uno SRAM left after the compiler-reported static use, the free heap at $n=2$ on the ESP8266, and on the ESP32 either the total free heap or the largest free block (one contiguous allocation is needed).

\textbf{Time models.} M1: $t=c\cdot\text{work}(n)$ with one constant per board. M1c: one constant per board and circuit. M2: $t=(c_H N_H + c_{CNOT}N_{CNOT} + c_{CP}N_{CP})2^n + aG$ with non-negative least squares on relative error, four parameters per board shared by the circuits ($G$: gate calls); M2g is M2 fitted without \grover. Per series the largest three $n$ are held out and every other point with median $\ge$\,\ModelFloorUs\us is fitted unless it is weak or hash-flagged. The split and the floor were fixed before any test point was scored (by the authors' record; not verifiable from the commit history). M2 has the form of an exploratory model whose errors on \grover\ were already known to the authors: it was written down before this analysis was run but \emph{not defined blind}. A fit is called identifiable if it has at least (parameters $+2$) points, full column rank, a column-scaled condition number $\le$ \CondCap\ and no coefficient at the non-negativity boundary; this rule and the weak, hash and effective-repetition rules were \emph{not} fixed before scoring.

\textbf{Crossover bands.} For each circuit and $n$ the observed minimum and maximum PC simulation time over \emph{all} Phase~3 sessions on every board span the offload latency: wire $+$ residual (selected session) $+$ simulation (min or max). The crossover is the bracket of two consecutive measured $n$ between which the offload band drops below the Phase~2 local median and stays below at every larger measured $n$, once for the fast and once for the slow bound. No interpolation or extrapolation. Labels: \emph{within noise}: local and offload IQR overlap at an end of the bracket; \emph{weak}: a weak point lies in or beyond the bracket, or supplies the extreme PC simulation time. The band design was fixed after the data were seen.
